\documentclass[../DoAn.tex]{subfiles}
\begin{document}

Chương 3 đã trình bày quá trình xây dựng dữ liệu, kết thúc ở các tensor tín hiệu quán tính đã chuẩn hóa cùng dữ liệu tương tác màn hình thô được tổ chức theo người dùng. Chương 4 trình bày quá trình xây dựng mô hình trên cơ sở dữ liệu đó. Phần \ref{sec:backbone} mô tả 3 kiến trúc backbone cho nhánh quán tính và cách huấn luyện chúng để học không gian embedding. Phần \ref{sec:touchmodel} trình bày việc trích xuất đặc trưng và phân loại dữ liệu tương tác màn hình bằng Random Forest. Phần \ref{sec:decision} mô tả hàm quyết định nhận dạng tập mở dựa trên đối sánh embedding, cơ chế dung hợp đa phương thức và bước làm mịn theo thời gian. Phần \ref{sec:export_tflite} trình bày việc xuất mô hình sang định dạng TensorFlow Lite phục vụ triển khai trên thiết bị.

\section{Backbone cho dữ liệu quán tính}
\label{sec:backbone}

\subsection{Ba kiến trúc khảo sát}

Đề tài khảo sát 3 kiến trúc backbone để mã hóa cửa sổ tín hiệu quán tính kích thước $200 \times 9$ thành một vector embedding 128 chiều. Kiến trúc thứ nhất là CNN 1D, gồm 3 khối tích chập một chiều với số kênh tăng dần từ 9 lên 64, 128 và 128, mỗi khối kèm chuẩn hóa theo lô và hàm kích hoạt ReLU; 2 khối đầu có tầng gộp cực đại để rút gọn chiều thời gian, khối cuối dùng gộp trung bình thích ứng để thu về một vector embedding cố định. Kiến trúc thứ hai là CNN-LSTM, dùng 2 khối tích chập một chiều để trích xuất đặc trưng cục bộ và rút gọn chuỗi, sau đó đưa qua một lớp LSTM để học phụ thuộc thời gian dài hạn. Kiến trúc thứ ba là CNN-BiLSTM, thay LSTM một chiều bằng LSTM 2 chiều để mỗi bước thời gian tổng hợp thông tin từ cả 2 phía của chuỗi.

Cả 3 kiến trúc dùng chung một giao diện: nhận đầu vào là cửa sổ tín hiệu 9 kênh và trả về embedding 128 chiều, nhờ đó có thể thay thế lẫn nhau trong cùng một pipeline huấn luyện và đánh giá. Như đã nêu ở Chương 2, 2 biến thể có thành phần hồi tiếp là kiến trúc lai tuần tự gồm tích chập theo sau mạng hồi tiếp, khác với mô hình ConvLSTM đặt tích chập bên trong cổng của ô LSTM. Ba kiến trúc này được lựa chọn để khảo sát một cách có hệ thống vai trò của thành phần mô hình hóa thời gian: CNN 1D thuần đại diện cho hướng chỉ trích đặc trưng cục bộ bằng tích chập, CNN-LSTM bổ sung khả năng học phụ thuộc thời gian dài hạn theo một chiều, còn CNN-BiLSTM mở rộng sang 2 chiều. Việc so sánh 3 mức độ phức tạp tăng dần này cho phép đánh giá liệu thành phần hồi tiếp có thực sự cải thiện hiệu năng đủ để bù lại chi phí tăng thêm về tham số và độ trễ hay không.

\subsection{Huấn luyện học không gian embedding}

Mục tiêu của giai đoạn huấn luyện không phải là phân loại trực tiếp người dùng lúc vận hành, mà là học một không gian embedding trong đó các cửa sổ của cùng một người nằm gần nhau và khác biệt với người khác. Để đạt mục tiêu đó, mỗi backbone được gắn tạm một tầng phân loại tuyến tính và huấn luyện như bài toán phân loại nhiều người dùng trên dữ liệu của toàn bộ người tham gia trong tập huấn luyện, sử dụng hàm mất mát entropy chéo. Nhãn hoạt động không được đưa vào hàm mất mát; backbone học biểu diễn độc lập với hoạt động thông qua việc quan sát dữ liệu của cả 3 hoạt động. Sau khi huấn luyện hoàn tất, tầng phân loại được loại bỏ và chỉ bộ mã hóa sinh embedding được giữ lại để triển khai.

Quá trình huấn luyện sử dụng thuật toán tối ưu Adam với tốc độ học khởi tạo $10^{-3}$ và hệ số suy giảm trọng số $10^{-4}$, kích thước lô 256, huấn luyện tối đa 30 vòng lặp với cơ chế dừng sớm khi độ chính xác trên tập kiểm định không cải thiện sau 8 vòng liên tiếp. Gradient được cắt ngưỡng theo chuẩn để ổn định huấn luyện. Để tăng tính bền vững của biểu diễn, dữ liệu huấn luyện được tăng cường bằng cách thêm nhiễu Gauss nhỏ và co giãn biên độ trong phạm vi hẹp. Toàn bộ quá trình được thực hiện 2 lần độc lập tương ứng 2 cấu hình ngữ cảnh đi bộ và toàn bộ 3 hoạt động, tạo thành 2 bộ mô hình để so sánh ở Chương 5.

\subsection{So sánh quy mô mô hình}

Vì hệ thống hướng tới suy luận trực tiếp trên thiết bị, quy mô mô hình là một tiêu chí lựa chọn quan trọng bên cạnh độ chính xác. Bảng~\ref{tab:model_size} so sánh số tham số và dung lượng của bộ mã hóa 3 kiến trúc.

\begin{table}[H]
    \centering
    \caption{So sánh quy mô 3 bộ mã hóa backbone}
    \label{tab:model_size}
    \begin{tabular}{|l|c|c|}
    \hline
    \textbf{Kiến trúc} & \textbf{Số tham số} & \textbf{Dung lượng (tệp checkpoint)} \\
    \hline
    CNN 1D     & $\approx$ 77,6 nghìn  & $\approx$ 331 KB \\
    \hline
    CNN-LSTM   & $\approx$ 160,1 nghìn & $\approx$ 659 KB \\
    \hline
    CNN-BiLSTM & $\approx$ 127,4 nghìn & $\approx$ 528 KB \\
    \hline
    \end{tabular}
\end{table}

Như thể hiện trong Bảng~\ref{tab:model_size}, CNN 1D là kiến trúc gọn nhẹ nhất, với số tham số chỉ bằng khoảng một nửa so với 2 biến thể có thành phần hồi tiếp. Cùng với kết quả đánh giá ở Chương 5, ưu thế về quy mô và tốc độ suy luận này là cơ sở để lựa chọn CNN 1D làm kiến trúc triển khai trên thiết bị.

\subsection{Cơ sở lựa chọn độ phức tạp mô hình theo kích thước dữ liệu}

Việc lựa chọn độ phức tạp của bộ mã hóa cần được đặt trong tương quan với quy mô dữ liệu huấn luyện, đặc biệt với mục tiêu nhận dạng tập mở. Bộ dữ liệu gồm 26 người tham gia, mỗi người 6 phiên; sau khi giữ riêng 5 danh tính cho đánh giá tập mở và loại một số ít hồ sơ thiếu dữ liệu, còn khoảng 19 danh tính được dùng để huấn luyện bộ mã hóa. Mặc dù số cửa sổ tín hiệu lên tới hàng chục nghìn, các cửa sổ này được trượt chồng lấn nên những cửa sổ trong cùng một phiên gần như trùng lặp về thông tin; do đó số mẫu thực sự độc lập nhỏ hơn nhiều con số cửa sổ thô và xấp xỉ số khối danh tính--phiên vào khoảng tích của 19 danh tính với 6 phiên, tức trên dưới một trăm khối. Đối với khả năng tổng quát hóa sang người lạ, yếu tố quyết định là độ đa dạng về danh tính và phiên hành vi, chứ không phải số lượng cửa sổ thô; và cả 2 đại lượng này đều ở mức hạn chế.

Từ tương quan đó, một mô hình có dung lượng quá lớn so với lượng danh tính quan sát được sẽ có xu hướng khớp với chính tập danh tính đã thấy và tổng quát hóa kém cho người chưa biết. Vì vậy, thiết kế ưu tiên một bộ mã hóa gọn nhẹ thay vì các kiến trúc nhiều tham số. Định hướng này nhất quán với so sánh quy mô ở Bảng~\ref{tab:model_size}, trong đó CNN 1D chỉ dùng khoảng một nửa số tham số so với 2 biến thể hồi tiếp và được khẳng định bằng thực nghiệm ở Chương 5: ở kịch bản tập mở, CNN 1D cho tỉ lệ lỗi cân bằng thấp nhất, trong khi CNN-LSTM và CNN-BiLSTM tuy nhiều tham số hơn lại không cải thiện mà còn có dấu hiệu khớp với tập danh tính đã quan sát. Như vậy, độ phức tạp mô hình được lựa chọn tương xứng với kích thước dữ liệu hiện có.

\section{Trích xuất đặc trưng và phân loại dữ liệu tương tác}
\label{sec:touchmodel}

\subsection{Trích xuất đặc trưng động học}

Khác với dữ liệu quán tính được đưa trực tiếp vào mạng học sâu, dữ liệu tương tác màn hình được chuyển thành các vector đặc trưng động học mô tả hành vi chạm và cuộn trước khi phân loại. Nhóm thao tác chạm gồm số lần chạm cùng các thống kê về thời gian giữ, độ dịch chuyển và khoảng thời gian giữa các lần chạm. Nhóm thao tác cuộn gồm số lần cuộn cùng các thống kê về thời lượng, độ dài và độ thẳng quỹ đạo, vận tốc trung bình và cực đại, hướng cuộn và tỉ lệ các hướng.

Với mỗi nhóm, các đặc trưng thống kê được tính dưới dạng trung bình, độ lệch chuẩn và các phân vị nhằm mô tả cả xu hướng trung tâm lẫn mức độ biến thiên của hành vi. Như đã nêu ở Chương 2, áp lực và diện tích tiếp xúc không được đưa vào tập đặc trưng vì trên phần lớn thiết bị thu thập các giá trị này gần như không đổi nên ít mang thông tin phân biệt.

\subsection{Phân loại bằng Random Forest}

Vector đặc trưng 48 chiều của nhánh touch được phân loại bằng thuật toán Random Forest theo chiến lược owner-vs-pool: với mỗi chủ sở hữu, một bộ phân loại nhị phân được huấn luyện để phân biệt các mẫu của chủ sở hữu với một tập mẫu nền lấy từ những người dùng khác. Bộ phân loại sử dụng 200 cây quyết định, số đặc trưng xét tại mỗi nút lấy theo căn bậc 2 tổng số đặc trưng và trọng số lớp được cân bằng để xử lý chênh lệch giữa số mẫu của chủ sở hữu và tập nền. Điểm tin cậy của một phiên được lấy bằng trung bình xác suất dự đoán trên các cửa sổ con của phiên đó. Random Forest được lựa chọn cho nhánh touch thay vì mạng học sâu vì lượng dữ liệu tương tác nhỏ hơn nhiều so với dữ liệu quán tính và một mô hình đơn giản giúp hạn chế quá khớp trên lượng dữ liệu hạn chế này.

\section{Hàm quyết định và dung hợp}
\label{sec:decision}

\subsection{Hàm quyết định nhận dạng tập mở}

Hàm quyết định của nhánh quán tính được thiết kế theo hướng đối sánh mẫu trong không gian embedding thay vì phân loại trực tiếp, nhằm phù hợp với kịch bản nhận dạng tập mở và cho phép đăng ký người dùng mới mà không cần huấn luyện lại mô hình. Khi đăng ký, một số cửa sổ của chủ sở hữu được đưa qua bộ mã hóa để sinh embedding, từ đó chọn ra một tập gồm các embedding đại diện gọi là anchor, ký hiệu $\{a_1, a_2, \dots, a_K\}$. Khi xác thực, embedding $z$ của cửa sổ đầu vào được so với từng anchor bằng độ tương đồng cosine

\begin{equation}
    \operatorname{sim}(z, a_k) = \frac{z \cdot a_k}{\lVert z \rVert \, \lVert a_k \rVert},
    \label{eq:cosine}
\end{equation}

và độ tương đồng tổng hợp được lấy trung bình trên toàn bộ anchor,

\begin{equation}
    \bar{s}(z) = \frac{1}{K} \sum_{k=1}^{K} \operatorname{sim}(z, a_k).
    \label{eq:cosine_mean}
\end{equation}

Giá trị tương đồng trung bình này được đưa qua một hàm sigmoid có hệ số co giãn $\gamma$ và độ dịch $\beta$ để chuyển thành điểm tin cậy trong khoảng đơn vị,

\begin{equation}
    s_\text{quán tính}(z) = \sigma\!\big(\gamma \, (\bar{s}(z) - \beta)\big),
    \qquad \sigma(u) = \frac{1}{1 + e^{-u}},
    \label{eq:sigmoid}
\end{equation}

trong đó hệ số $\gamma$ điều chỉnh độ dốc và $\beta$ là ngưỡng tương đồng tham chiếu. Người dùng lạ chưa từng đăng ký sẽ cho embedding ít tương đồng với các anchor của chủ sở hữu, dẫn đến điểm tin cậy thấp và bị từ chối khi điểm nằm dưới ngưỡng quyết định.

\subsection{Dung hợp đa phương thức}

Để khảo sát đóng góp của dữ liệu tương tác, đề tài xây dựng cấu hình dung hợp ở mức điểm số, trong đó điểm tin cậy của 2 nhánh được kết hợp bằng trung bình có trọng số

\begin{equation}
    s = w \cdot s_\text{quán tính} + (1 - w) \cdot s_\text{touch},
    \qquad w \in [0, 1],
    \label{eq:fusion}
\end{equation}

với $s_\text{quán tính}$ và $s_\text{touch}$ lần lượt là điểm tin cậy của nhánh quán tính và nhánh tương tác. Trọng số $w$ được tìm kiếm trên tập kiểm định bằng cách duyệt lưới các giá trị trong khoảng đơn vị để cực đại hóa diện tích dưới đường cong ROC; khi nhiều giá trị cho cùng kết quả, trọng số gần mức cân bằng được ưu tiên để tránh thoái hóa về một nhánh duy nhất. Như sẽ phân tích ở Chương 5, cấu hình dung hợp này không cải thiện hiệu năng trong kịch bản tập mở, nên đóng vai trò đối chứng còn cấu hình triển khai cuối cùng chỉ sử dụng nhánh quán tính.

\subsection{Làm mịn điểm theo thời gian}

Một quyết định xác thực dựa trên một cửa sổ đơn lẻ dễ bị ảnh hưởng bởi nhiễu nhất thời. Do đó, điểm tin cậy được làm mịn theo thời gian bằng trung bình động trọng số mũ trên một cửa sổ gồm các điểm gần nhất. Gọi $s_t$ là điểm tin cậy thô tại bước $t$ và $\hat{s}_t$ là điểm sau khi làm mịn, cơ chế làm mịn được biểu diễn bằng công thức truy hồi

\begin{equation}
    \hat{s}_t = \alpha \, s_t + (1 - \alpha) \, \hat{s}_{t-1},
    \qquad \alpha \in (0, 1),
    \label{eq:ewma}
\end{equation}

trong đó hệ số làm mượt $\alpha$ quyết định mức độ ưu tiên điểm mới so với lịch sử: $\alpha$ càng lớn thì điểm mới nhất càng có trọng số cao và các điểm cũ giảm dần theo lũy thừa của $(1 - \alpha)$. Cửa sổ làm mịn được giới hạn ở vài điểm gần nhất để cân bằng giữa 2 yêu cầu đối nghịch: cửa sổ quá ngắn khiến điểm số vẫn nhạy với nhiễu, còn cửa sổ quá dài làm hệ thống phản ứng chậm khi người dùng thực sự thay đổi. Hệ số $\alpha$ được đặt thiên về điểm gần nhất để hệ thống vừa giữ được tính ổn định, vừa ưu tiên thông tin mới và phát hiện kịp thời thay đổi hành vi. Điểm sau khi làm mịn được so với ngưỡng quyết định để đưa ra kết luận chấp nhận hay từ chối; ngưỡng được xác định tại điểm cân bằng sai số trên tập đánh giá, như trình bày ở Chương 5. Riêng bản triển khai trên thiết bị còn hiệu chuẩn ngưỡng này riêng cho từng chủ sở hữu ngay khi đăng ký, thay cho một mức chung, như trình bày ở Chương 6.

\section{Xuất mô hình sang TensorFlow Lite}
\label{sec:export_tflite}

Để phục vụ suy luận trực tiếp trên thiết bị, bộ mã hóa CNN 1D sau huấn luyện được chuyển đổi từ định dạng PyTorch sang định dạng TensorFlow Lite. Mô hình được giữ ở độ chính xác dấu phẩy động 32 bit nhằm bảo toàn độ chính xác của embedding, vốn là cơ sở cho phép đối sánh cosine trong hàm quyết định. Kèm theo tệp mô hình, các tham số cần thiết cho suy luận như tập anchor của chủ sở hữu và bộ chuẩn hóa đặc trưng cũng được đóng gói để tích hợp vào ứng dụng xác thực.

Việc giữ độ chính xác dấu phẩy động đầy đủ thay vì lượng tử hóa xuống số nguyên là một đánh đổi có chủ đích: do bộ mã hóa đã gọn nhẹ, dung lượng và độ trễ ở mức chấp nhận được trên thiết bị, trong khi việc lượng tử hóa có thể làm sai lệch embedding và ảnh hưởng đến độ tương đồng cosine. Việc khảo sát lượng tử hóa nhằm giảm thêm dung lượng mà vẫn giữ chất lượng đối sánh được xem là một hướng cải thiện trong tương lai.

Chương này đã trình bày toàn bộ quá trình xây dựng mô hình, từ 3 kiến trúc backbone cho dữ liệu quán tính, việc trích xuất đặc trưng và phân loại dữ liệu tương tác, đến hàm quyết định nhận dạng tập mở, cơ chế dung hợp và bước xuất mô hình sang TensorFlow Lite. Trên cơ sở các mô hình đã xây dựng, Chương 5 trình bày thiết kế thực nghiệm và kết quả đánh giá, so sánh 3 kiến trúc backbone trên 2 kịch bản nhận dạng tập đóng và tập mở, từ đó rút ra cấu hình phù hợp nhất để hiện thực hóa trên ứng dụng Android ở Chương 6.

\end{document}